\section{Erd\H{o}s Problem \#194}
\subsection*{1. FORMAL RESTATEMENT}
For $k\ge3$, must every linear order $\prec$ on $\mathbb R$ have a nonconstant arithmetic progression $x,x+d,\ldots,x+(k-1)d$ that is monotone in $\prec$?
\subsection*{2. LITERATURE AND APPROACH}
Ardal, Brown and Jungi\'c proved a negative answer already for $k=3$. We give the rational order explicitly by binary digits in the two-adic direction and then use their Hamel-basis lifting idea. All order and progression checks are included. The construction uses the axiom of choice through the existence of a basis of $\mathbb R$ over $\mathbb Q$; it is not a claim of a new result.
\subsection*{3. AN EXPLICIT ORDER ON THE RATIONALS}
Every rational $q$ has a unique two-adic digit sequence
\[
 (\epsilon_j(q))_{j\in\mathbb Z},\qquad \epsilon_j(q)\in\{0,1\},
 \qquad\epsilon_j(q)=0\text{ for all sufficiently negative }j.
\]
This notation need not presume a theorem about complete fields. Choose $J$ with $2^{-J}q$ having odd denominator. For every $t\ge1$, reduce that rational modulo $2^t$ by inverting its odd denominator; its compatible residues determine the digits at $J,J+1,\ldots$. All lower digits are zero. Changing $J$ to a smaller permissible value leaves the same digits, by compatibility.

Two different rationals first differ at exponent $v_2(q-r)$, a finite integer. Define $q\prec_{\mathbb Q}r$ when at their first differing exponent the digit of $q$ is zero and that of $r$ is one. This is a strict linear order. For transitivity, if $q\prec r\prec s$, examine the smaller of the first differing exponents for $(q,r)$ and $(r,s)$. If unequal, the earlier comparison also determines $(q,s)$; if equal, the middle digit would have to be both one and zero, which is impossible. Trichotomy and irreflexivity are immediate.

For $d\ne0$, put $v_2(d)=t$. The rationals $x$ and $x+2d$ have identical digits through exponent $t$, whereas $x+d$ agrees with them below $t$ and has the opposite digit at $t$. Thus $x+d$ lies either before both endpoints or after both endpoints. No nonconstant three-term progression is monotone in this order.
\subsection*{4. LIFTING TO THE REALS}
Choose a Hamel basis $B$ of $\mathbb R$ over $\mathbb Q$ and linearly order its elements by their usual real values. Write $x=\sum_{b\in B}x_b b$, with rational coefficients and finite support. For distinct $x,y$, let $b_*$ be the largest element in the finite set $\{b:x_b\ne y_b\}$. Declare
\[
 x\prec y\quad\Longleftrightarrow\quad x_{b_*}\prec_{\mathbb Q}y_{b_*}.
\]
This finite-support reverse lexicographic rule is a linear order. To verify transitivity for three real numbers, restrict to the finite union of their supports and apply the ordinary lexicographic comparison of a finite product of linear orders, reading basis coordinates from largest to smallest.

For $x,x+d,x+2d$ with $d\ne0$, choose the largest basis element in the support of $d$. All higher coordinates of the three numbers agree. At the chosen coordinate their coefficients form a nonconstant rational progression. Its midpoint lies outside its endpoints in $\prec_{\mathbb Q}$, so the same is true in $\prec$. No monotone three-term progression exists, and a monotone longer progression would contain one among its first three terms.
\subsection*{5. VERIFICATION AND FINAL}
\textbf{SOLVED: negative for every $k\ge3$, in the usual set theory with choice.}
The digit sequences are not ordinary real binary expansions, and are never summed as convergent real series. Negative rationals and odd denominators are covered by modular inversion. The Hamel order is an arbitrary linear order, not required to respect addition or the usual topology. No unproved combinatorial input remains.
\subsection*{6. SOURCES}
Current statement: \texttt{https://www.erdosproblems.com/194}. H. Ardal, T. Brown and V. Jungi\'c, \emph{Chaotic orderings of the rationals and reals}, American Mathematical Monthly 118 (2011), 921--925, \texttt{https://www.sfu.ca/\~{}vjungic/tbrown/tom-3.pdf}. Their rational construction uses a compactness step; the digit construction above supplies that rational order directly.
\section{COMPLETION ESTIMATE}
\textbf{COMPLETION: 100\%}
